% Options for packages loaded elsewhere
% Options for packages loaded elsewhere
\PassOptionsToPackage{unicode}{hyperref}
\PassOptionsToPackage{hyphens}{url}
\PassOptionsToPackage{dvipsnames,svgnames,x11names}{xcolor}
%
\documentclass[
  letterpaper,
  DIV=11,
  numbers=noendperiod]{scrartcl}
\usepackage{xcolor}
\usepackage{amsmath,amssymb}
\setcounter{secnumdepth}{-\maxdimen} % remove section numbering
\usepackage{iftex}
\ifPDFTeX
  \usepackage[T1]{fontenc}
  \usepackage[utf8]{inputenc}
  \usepackage{textcomp} % provide euro and other symbols
\else % if luatex or xetex
  \usepackage{unicode-math} % this also loads fontspec
  \defaultfontfeatures{Scale=MatchLowercase}
  \defaultfontfeatures[\rmfamily]{Ligatures=TeX,Scale=1}
\fi
\usepackage{lmodern}
\ifPDFTeX\else
  % xetex/luatex font selection
\fi
% Use upquote if available, for straight quotes in verbatim environments
\IfFileExists{upquote.sty}{\usepackage{upquote}}{}
\IfFileExists{microtype.sty}{% use microtype if available
  \usepackage[]{microtype}
  \UseMicrotypeSet[protrusion]{basicmath} % disable protrusion for tt fonts
}{}
\makeatletter
\@ifundefined{KOMAClassName}{% if non-KOMA class
  \IfFileExists{parskip.sty}{%
    \usepackage{parskip}
  }{% else
    \setlength{\parindent}{0pt}
    \setlength{\parskip}{6pt plus 2pt minus 1pt}}
}{% if KOMA class
  \KOMAoptions{parskip=half}}
\makeatother
% Make \paragraph and \subparagraph free-standing
\makeatletter
\ifx\paragraph\undefined\else
  \let\oldparagraph\paragraph
  \renewcommand{\paragraph}{
    \@ifstar
      \xxxParagraphStar
      \xxxParagraphNoStar
  }
  \newcommand{\xxxParagraphStar}[1]{\oldparagraph*{#1}\mbox{}}
  \newcommand{\xxxParagraphNoStar}[1]{\oldparagraph{#1}\mbox{}}
\fi
\ifx\subparagraph\undefined\else
  \let\oldsubparagraph\subparagraph
  \renewcommand{\subparagraph}{
    \@ifstar
      \xxxSubParagraphStar
      \xxxSubParagraphNoStar
  }
  \newcommand{\xxxSubParagraphStar}[1]{\oldsubparagraph*{#1}\mbox{}}
  \newcommand{\xxxSubParagraphNoStar}[1]{\oldsubparagraph{#1}\mbox{}}
\fi
\makeatother

\usepackage{color}
\usepackage{fancyvrb}
\newcommand{\VerbBar}{|}
\newcommand{\VERB}{\Verb[commandchars=\\\{\}]}
\DefineVerbatimEnvironment{Highlighting}{Verbatim}{commandchars=\\\{\}}
% Add ',fontsize=\small' for more characters per line
\usepackage{framed}
\definecolor{shadecolor}{RGB}{241,243,245}
\newenvironment{Shaded}{\begin{snugshade}}{\end{snugshade}}
\newcommand{\AlertTok}[1]{\textcolor[rgb]{0.68,0.00,0.00}{#1}}
\newcommand{\AnnotationTok}[1]{\textcolor[rgb]{0.37,0.37,0.37}{#1}}
\newcommand{\AttributeTok}[1]{\textcolor[rgb]{0.40,0.45,0.13}{#1}}
\newcommand{\BaseNTok}[1]{\textcolor[rgb]{0.68,0.00,0.00}{#1}}
\newcommand{\BuiltInTok}[1]{\textcolor[rgb]{0.00,0.23,0.31}{#1}}
\newcommand{\CharTok}[1]{\textcolor[rgb]{0.13,0.47,0.30}{#1}}
\newcommand{\CommentTok}[1]{\textcolor[rgb]{0.37,0.37,0.37}{#1}}
\newcommand{\CommentVarTok}[1]{\textcolor[rgb]{0.37,0.37,0.37}{\textit{#1}}}
\newcommand{\ConstantTok}[1]{\textcolor[rgb]{0.56,0.35,0.01}{#1}}
\newcommand{\ControlFlowTok}[1]{\textcolor[rgb]{0.00,0.23,0.31}{\textbf{#1}}}
\newcommand{\DataTypeTok}[1]{\textcolor[rgb]{0.68,0.00,0.00}{#1}}
\newcommand{\DecValTok}[1]{\textcolor[rgb]{0.68,0.00,0.00}{#1}}
\newcommand{\DocumentationTok}[1]{\textcolor[rgb]{0.37,0.37,0.37}{\textit{#1}}}
\newcommand{\ErrorTok}[1]{\textcolor[rgb]{0.68,0.00,0.00}{#1}}
\newcommand{\ExtensionTok}[1]{\textcolor[rgb]{0.00,0.23,0.31}{#1}}
\newcommand{\FloatTok}[1]{\textcolor[rgb]{0.68,0.00,0.00}{#1}}
\newcommand{\FunctionTok}[1]{\textcolor[rgb]{0.28,0.35,0.67}{#1}}
\newcommand{\ImportTok}[1]{\textcolor[rgb]{0.00,0.46,0.62}{#1}}
\newcommand{\InformationTok}[1]{\textcolor[rgb]{0.37,0.37,0.37}{#1}}
\newcommand{\KeywordTok}[1]{\textcolor[rgb]{0.00,0.23,0.31}{\textbf{#1}}}
\newcommand{\NormalTok}[1]{\textcolor[rgb]{0.00,0.23,0.31}{#1}}
\newcommand{\OperatorTok}[1]{\textcolor[rgb]{0.37,0.37,0.37}{#1}}
\newcommand{\OtherTok}[1]{\textcolor[rgb]{0.00,0.23,0.31}{#1}}
\newcommand{\PreprocessorTok}[1]{\textcolor[rgb]{0.68,0.00,0.00}{#1}}
\newcommand{\RegionMarkerTok}[1]{\textcolor[rgb]{0.00,0.23,0.31}{#1}}
\newcommand{\SpecialCharTok}[1]{\textcolor[rgb]{0.37,0.37,0.37}{#1}}
\newcommand{\SpecialStringTok}[1]{\textcolor[rgb]{0.13,0.47,0.30}{#1}}
\newcommand{\StringTok}[1]{\textcolor[rgb]{0.13,0.47,0.30}{#1}}
\newcommand{\VariableTok}[1]{\textcolor[rgb]{0.07,0.07,0.07}{#1}}
\newcommand{\VerbatimStringTok}[1]{\textcolor[rgb]{0.13,0.47,0.30}{#1}}
\newcommand{\WarningTok}[1]{\textcolor[rgb]{0.37,0.37,0.37}{\textit{#1}}}

\usepackage{longtable,booktabs,array}
\usepackage{calc} % for calculating minipage widths
% Correct order of tables after \paragraph or \subparagraph
\usepackage{etoolbox}
\makeatletter
\patchcmd\longtable{\par}{\if@noskipsec\mbox{}\fi\par}{}{}
\makeatother
% Allow footnotes in longtable head/foot
\IfFileExists{footnotehyper.sty}{\usepackage{footnotehyper}}{\usepackage{footnote}}
\makesavenoteenv{longtable}
\usepackage{graphicx}
\makeatletter
\newsavebox\pandoc@box
\newcommand*\pandocbounded[1]{% scales image to fit in text height/width
  \sbox\pandoc@box{#1}%
  \Gscale@div\@tempa{\textheight}{\dimexpr\ht\pandoc@box+\dp\pandoc@box\relax}%
  \Gscale@div\@tempb{\linewidth}{\wd\pandoc@box}%
  \ifdim\@tempb\p@<\@tempa\p@\let\@tempa\@tempb\fi% select the smaller of both
  \ifdim\@tempa\p@<\p@\scalebox{\@tempa}{\usebox\pandoc@box}%
  \else\usebox{\pandoc@box}%
  \fi%
}
% Set default figure placement to htbp
\def\fps@figure{htbp}
\makeatother





\setlength{\emergencystretch}{3em} % prevent overfull lines

\providecommand{\tightlist}{%
  \setlength{\itemsep}{0pt}\setlength{\parskip}{0pt}}



 


\KOMAoption{captions}{tableheading}
\makeatletter
\@ifpackageloaded{caption}{}{\usepackage{caption}}
\AtBeginDocument{%
\ifdefined\contentsname
  \renewcommand*\contentsname{Table of contents}
\else
  \newcommand\contentsname{Table of contents}
\fi
\ifdefined\listfigurename
  \renewcommand*\listfigurename{List of Figures}
\else
  \newcommand\listfigurename{List of Figures}
\fi
\ifdefined\listtablename
  \renewcommand*\listtablename{List of Tables}
\else
  \newcommand\listtablename{List of Tables}
\fi
\ifdefined\figurename
  \renewcommand*\figurename{Figure}
\else
  \newcommand\figurename{Figure}
\fi
\ifdefined\tablename
  \renewcommand*\tablename{Table}
\else
  \newcommand\tablename{Table}
\fi
}
\@ifpackageloaded{float}{}{\usepackage{float}}
\floatstyle{ruled}
\@ifundefined{c@chapter}{\newfloat{codelisting}{h}{lop}}{\newfloat{codelisting}{h}{lop}[chapter]}
\floatname{codelisting}{Listing}
\newcommand*\listoflistings{\listof{codelisting}{List of Listings}}
\makeatother
\makeatletter
\makeatother
\makeatletter
\@ifpackageloaded{caption}{}{\usepackage{caption}}
\@ifpackageloaded{subcaption}{}{\usepackage{subcaption}}
\makeatother
\usepackage{bookmark}
\IfFileExists{xurl.sty}{\usepackage{xurl}}{} % add URL line breaks if available
\urlstyle{same}
\hypersetup{
  pdftitle={Class 6 R: functions},
  pdfauthor={David Majeed (PID:A17885958)},
  colorlinks=true,
  linkcolor={blue},
  filecolor={Maroon},
  citecolor={Blue},
  urlcolor={Blue},
  pdfcreator={LaTeX via pandoc}}


\title{Class 6 R: functions}
\author{David Majeed (PID:A17885958)}
\date{}
\begin{document}
\maketitle


\subsection{Background}\label{background}

All functions in R have 3 things:

-a \emph{name} (we choose)

-input \emph{arguments} (one or more comma-separated inputs that go
inside the brackets)

-the \emph{body} (the lines of R code that do the work of the function)

\subsection{First Function}\label{first-function}

Here we will create a function to add some numbers. Called
\texttt{add()}

Can be \textbf{``required'' or ``optional''}, the later will have
fall-back default values?

Q.1A

\begin{Shaded}
\begin{Highlighting}[]
\NormalTok{add}\OtherTok{\textless{}{-}}\ControlFlowTok{function}\NormalTok{(x,}\AttributeTok{y=}\DecValTok{1000}\NormalTok{)\{}
\NormalTok{  x}\SpecialCharTok{+}\NormalTok{y}
\NormalTok{\}}

\CommentTok{\#add is the name, function() is the function we use to make new function, the (x,y=1000) equal the inputs we are looking for, the =1000 is the defualt option if the user doesn\textquotesingle{}t inpiut anyting, \{\} is what the function actually does}
\end{Highlighting}
\end{Shaded}

Here we have the name: add, the inputs: x and y and the body of code:
x+y

Can we use it?

\begin{Shaded}
\begin{Highlighting}[]
\FunctionTok{add}\NormalTok{(}\DecValTok{10}\NormalTok{,}\DecValTok{100}\NormalTok{)}
\end{Highlighting}
\end{Shaded}

\begin{verbatim}
[1] 110
\end{verbatim}

\begin{Shaded}
\begin{Highlighting}[]
\FunctionTok{add}\NormalTok{(}\DecValTok{10}\NormalTok{)}
\end{Highlighting}
\end{Shaded}

\begin{verbatim}
[1] 1010
\end{verbatim}

\begin{Shaded}
\begin{Highlighting}[]
\CommentTok{\#we are testing our function}
\end{Highlighting}
\end{Shaded}

Q1B. Lets modify the add function

This allows there to be a default option in case the user forgot to add
one of the inputs

\begin{Shaded}
\begin{Highlighting}[]
\NormalTok{add}\OtherTok{\textless{}{-}}\ControlFlowTok{function}\NormalTok{(x,}\AttributeTok{y=}\DecValTok{0}\NormalTok{)\{}
  \FunctionTok{sum}\NormalTok{(x)}\SpecialCharTok{+}\FunctionTok{sum}\NormalTok{(y)}
\NormalTok{\}}
\CommentTok{\#modifying it to have the body be using the sum() function, that way it can add vectors as seen bellow}
\FunctionTok{add}\NormalTok{(}\DecValTok{4}\NormalTok{,}\DecValTok{7}\NormalTok{)}
\end{Highlighting}
\end{Shaded}

\begin{verbatim}
[1] 11
\end{verbatim}

\begin{Shaded}
\begin{Highlighting}[]
\FunctionTok{add}\NormalTok{(}\FunctionTok{c}\NormalTok{(}\DecValTok{4}\NormalTok{,}\DecValTok{7}\NormalTok{,}\DecValTok{10}\NormalTok{))}
\end{Highlighting}
\end{Shaded}

\begin{verbatim}
[1] 21
\end{verbatim}

Q.1C Lets add more inputs

This allows the user to add as many terms as they would like by using
\ldots{}

\begin{Shaded}
\begin{Highlighting}[]
\NormalTok{add}\OtherTok{\textless{}{-}}\ControlFlowTok{function}\NormalTok{(x,}\AttributeTok{y=}\DecValTok{0}\NormalTok{, ...)\{}
  \FunctionTok{sum}\NormalTok{(x)}\SpecialCharTok{+}\FunctionTok{sum}\NormalTok{(y)}\SpecialCharTok{+}\FunctionTok{sum}\NormalTok{(...)}
  
\NormalTok{\}}
\CommentTok{\#the ... allow us to have unlimited inputs, that way we avoid errors if the user doesn\textquotesingle{}t input a variable or if they have more inputs than we coded for}
\FunctionTok{add}\NormalTok{(}\DecValTok{1}\NormalTok{,}\DecValTok{2}\NormalTok{,}\DecValTok{3}\NormalTok{,}\DecValTok{4}\NormalTok{,}\DecValTok{5}\NormalTok{,}\DecValTok{7}\NormalTok{,}\DecValTok{7}\NormalTok{,}\DecValTok{7}\NormalTok{)}
\end{Highlighting}
\end{Shaded}

\begin{verbatim}
[1] 36
\end{verbatim}

\begin{Shaded}
\begin{Highlighting}[]
\FunctionTok{add}\NormalTok{(}\DecValTok{4}\NormalTok{, }\DecValTok{7}\NormalTok{)}
\end{Highlighting}
\end{Shaded}

\begin{verbatim}
[1] 11
\end{verbatim}

\begin{Shaded}
\begin{Highlighting}[]
\FunctionTok{add}\NormalTok{( }\FunctionTok{c}\NormalTok{(}\DecValTok{4}\NormalTok{,}\DecValTok{7}\NormalTok{,}\DecValTok{10}\NormalTok{))}
\end{Highlighting}
\end{Shaded}

\begin{verbatim}
[1] 21
\end{verbatim}

\begin{Shaded}
\begin{Highlighting}[]
\FunctionTok{add}\NormalTok{(}\DecValTok{1}\NormalTok{, }\DecValTok{2}\NormalTok{, }\DecValTok{3}\NormalTok{, }\SpecialCharTok{{-}}\DecValTok{4}\NormalTok{)}
\end{Highlighting}
\end{Shaded}

\begin{verbatim}
[1] 2
\end{verbatim}

We can \textbf{return} value from a function (rather than the default
last line) with \texttt{return()}

This helps us go back and just return the line we want

\begin{Shaded}
\begin{Highlighting}[]
\NormalTok{add}\OtherTok{\textless{}{-}}\ControlFlowTok{function}\NormalTok{(x,}\AttributeTok{y=}\DecValTok{0}\NormalTok{, ...)\{}
  \FunctionTok{return}\NormalTok{(}\FunctionTok{sum}\NormalTok{(x)}\SpecialCharTok{+}\FunctionTok{sum}\NormalTok{(y)}\SpecialCharTok{+}\FunctionTok{sum}\NormalTok{(...))}
  \FunctionTok{cat}\NormalTok{(}\StringTok{"is it break time?"}\NormalTok{)}
\NormalTok{\}}
\CommentTok{\#return() prevents the code from going further}
\FunctionTok{add}\NormalTok{(}\DecValTok{2}\NormalTok{,}\DecValTok{2}\NormalTok{)}
\end{Highlighting}
\end{Shaded}

\begin{verbatim}
[1] 4
\end{verbatim}

As we see, the text isn't printed

\subsection{A generate\_dna() function}\label{a-generate_dna-function}

We can use the \texttt{sample()} function

This allows us to randomly choose a value from a vector or set, the size
allows us to choose how many we select

\begin{Shaded}
\begin{Highlighting}[]
\FunctionTok{sample}\NormalTok{(}\DecValTok{1}\SpecialCharTok{:}\DecValTok{5}\NormalTok{, }\AttributeTok{size=}\DecValTok{5}\NormalTok{)}
\end{Highlighting}
\end{Shaded}

\begin{verbatim}
[1] 5 1 2 3 4
\end{verbatim}

Q.2A We can use this to make a random nucleotide sequence if we draw
from ``A'', ``C'', ``G'' and ``T''

\begin{Shaded}
\begin{Highlighting}[]
\CommentTok{\#sample(c("A","C","G","T"), 1, replace = TRUE)}

\NormalTok{generate\_dna }\OtherTok{\textless{}{-}} \ControlFlowTok{function}\NormalTok{(len)\{}
  \FunctionTok{sample}\NormalTok{(}\FunctionTok{c}\NormalTok{(}\StringTok{"A"}\NormalTok{,}\StringTok{"C"}\NormalTok{,}\StringTok{"G"}\NormalTok{,}\StringTok{"T"}\NormalTok{), len, }\AttributeTok{replace =} \ConstantTok{TRUE}\NormalTok{)}
\NormalTok{\}}

\CommentTok{\#we created a new function named generate\_dna that takes len as the input, and in return it gives a sample out of the avilable letters, for len \# and we can replace letters after chosing them}
\FunctionTok{generate\_dna}\NormalTok{(}\AttributeTok{len=}\DecValTok{6}\NormalTok{)}
\end{Highlighting}
\end{Shaded}

\begin{verbatim}
[1] "A" "G" "T" "G" "T" "A"
\end{verbatim}

Instead of using numbers 1:5, we can use nucleotides. The number of
nucleotides is selected from len input, and we are able to use more than
4 nucleotides by repeating them when replace=T

Q.2B

\begin{Shaded}
\begin{Highlighting}[]
\NormalTok{generate\_dna }\OtherTok{\textless{}{-}} \ControlFlowTok{function}\NormalTok{(len,single.element)\{}
\NormalTok{ans }\OtherTok{\textless{}{-}} \FunctionTok{sample}\NormalTok{(}\FunctionTok{c}\NormalTok{(}\StringTok{"A"}\NormalTok{,}\StringTok{"C"}\NormalTok{,}\StringTok{"G"}\NormalTok{,}\StringTok{"T"}\NormalTok{), len, }\AttributeTok{replace =} \ConstantTok{TRUE}\NormalTok{)}
    \ControlFlowTok{if}\NormalTok{ (single.element)\{}
     \FunctionTok{cat}\NormalTok{(}\FunctionTok{sample}\NormalTok{(}\FunctionTok{c}\NormalTok{(}\StringTok{"A"}\NormalTok{,}\StringTok{"C"}\NormalTok{,}\StringTok{"G"}\NormalTok{,}\StringTok{"T"}\NormalTok{), len, }\AttributeTok{replace =} \ConstantTok{TRUE}\NormalTok{), }\AttributeTok{sep =} \StringTok{""}\NormalTok{)\}}
\ControlFlowTok{else}\NormalTok{ \{}\FunctionTok{return}\NormalTok{(ans)\}}
\NormalTok{\}}

\CommentTok{\#if and else allow us to code for if the user wants to do an option or do the other one}
\CommentTok{\#sep removes the text to make a single string}

\FunctionTok{generate\_dna}\NormalTok{(}\AttributeTok{len=}\DecValTok{2}\NormalTok{, }\AttributeTok{single.element=}\NormalTok{T)}
\end{Highlighting}
\end{Shaded}

\begin{verbatim}
GA
\end{verbatim}

Building off the previous function, now we use \texttt{if} and
\texttt{else}. These allow us to control what happens when a certain
condition happens such as if we want to have a single character element
which then goes through \texttt{cat} to turn it into a string of
characters we want

Q2.C

\subsection{Fasta Format}\label{fasta-format}

\begin{Shaded}
\begin{Highlighting}[]
\DocumentationTok{\#\#Format as FASTA}

\NormalTok{generate\_dna }\OtherTok{\textless{}{-}} \ControlFlowTok{function}\NormalTok{(len,}\AttributeTok{single.element=}\NormalTok{T)\{}
  \FunctionTok{cat}\NormalTok{(}\FunctionTok{paste}\NormalTok{(}\StringTok{"\textgreater{}len"}\NormalTok{, len, }\StringTok{"}\SpecialCharTok{\textbackslash{}n}\StringTok{"}\NormalTok{, }\AttributeTok{sep=}\StringTok{""}\NormalTok{))}
\NormalTok{ans }\OtherTok{\textless{}{-}} \FunctionTok{sample}\NormalTok{(}\FunctionTok{c}\NormalTok{(}\StringTok{"A"}\NormalTok{,}\StringTok{"C"}\NormalTok{,}\StringTok{"G"}\NormalTok{,}\StringTok{"T"}\NormalTok{), len, }\AttributeTok{replace =} \ConstantTok{TRUE}\NormalTok{)}
    \ControlFlowTok{if}\NormalTok{ (single.element)\{}
     \FunctionTok{cat}\NormalTok{(}\FunctionTok{sample}\NormalTok{(}\FunctionTok{c}\NormalTok{(}\StringTok{"A"}\NormalTok{,}\StringTok{"C"}\NormalTok{,}\StringTok{"G"}\NormalTok{,}\StringTok{"T"}\NormalTok{), len, }\AttributeTok{replace =} \ConstantTok{TRUE}\NormalTok{), }\AttributeTok{sep =} \StringTok{""}\NormalTok{)\}}
\ControlFlowTok{else}\NormalTok{ \{}\FunctionTok{return}\NormalTok{(ans)\}}
\FunctionTok{cat}\NormalTok{(}\StringTok{"}\SpecialCharTok{\textbackslash{}n}\StringTok{"}\NormalTok{)}

\NormalTok{\}}

\CommentTok{\#cat controls caracters, such as pasting the format of "\textgreater{}len", then adding the number from length, then "\textbackslash{}n" goes down a line}

\FunctionTok{generate\_dna}\NormalTok{(}\AttributeTok{len=}\DecValTok{2}\NormalTok{, }\AttributeTok{single.element=}\NormalTok{T)}
\end{Highlighting}
\end{Shaded}

\begin{verbatim}
>len2
TT
\end{verbatim}

Now putting it all together, the above code use \texttt{cat} and
\texttt{paste} to now format it properly

Q.3

\subsection{generate\_protein()
function}\label{generate_protein-function}

\begin{Shaded}
\begin{Highlighting}[]
\NormalTok{generate\_protein }\OtherTok{\textless{}{-}} \ControlFlowTok{function}\NormalTok{(len,}\AttributeTok{single.element=}\NormalTok{T)\{}
\FunctionTok{cat}\NormalTok{(}\FunctionTok{paste}\NormalTok{(}\StringTok{"\textgreater{}len"}\NormalTok{, len, }\StringTok{"}\SpecialCharTok{\textbackslash{}n}\StringTok{"}\NormalTok{, }\AttributeTok{sep=}\StringTok{""}\NormalTok{))}
\NormalTok{ans }\OtherTok{\textless{}{-}} \FunctionTok{sample}\NormalTok{(}\FunctionTok{c}\NormalTok{(}\StringTok{"A"}\NormalTok{, }\StringTok{"R"}\NormalTok{, }\StringTok{"N"}\NormalTok{, }\StringTok{"D"}\NormalTok{, }\StringTok{"C"}\NormalTok{, }\StringTok{"E"}\NormalTok{, }\StringTok{"Q"}\NormalTok{, }\StringTok{"G"}\NormalTok{, }\StringTok{"H"}\NormalTok{, }\StringTok{"I"}\NormalTok{, }\StringTok{"L"}\NormalTok{, }\StringTok{"K"}\NormalTok{, }\StringTok{"M"}\NormalTok{, }\StringTok{"F"}\NormalTok{, }\StringTok{"P"}\NormalTok{, }\StringTok{"S"}\NormalTok{,}
\StringTok{"T"}\NormalTok{, }\StringTok{"W"}\NormalTok{, }\StringTok{"Y"}\NormalTok{, }\StringTok{"V"}\NormalTok{), len, }\AttributeTok{replace =} \ConstantTok{TRUE}\NormalTok{)}
    \ControlFlowTok{if}\NormalTok{ (single.element)\{}
     \FunctionTok{cat}\NormalTok{(}\FunctionTok{sample}\NormalTok{(}\FunctionTok{c}\NormalTok{(}\StringTok{"A"}\NormalTok{, }\StringTok{"R"}\NormalTok{, }\StringTok{"N"}\NormalTok{, }\StringTok{"D"}\NormalTok{, }\StringTok{"C"}\NormalTok{, }\StringTok{"E"}\NormalTok{, }\StringTok{"Q"}\NormalTok{, }\StringTok{"G"}\NormalTok{, }\StringTok{"H"}\NormalTok{, }\StringTok{"I"}\NormalTok{, }\StringTok{"L"}\NormalTok{, }\StringTok{"K"}\NormalTok{, }\StringTok{"M"}\NormalTok{, }\StringTok{"F"}\NormalTok{, }\StringTok{"P"}\NormalTok{, }\StringTok{"S"}\NormalTok{,}
\StringTok{"T"}\NormalTok{, }\StringTok{"W"}\NormalTok{, }\StringTok{"Y"}\NormalTok{, }\StringTok{"V"}\NormalTok{), len, }\AttributeTok{replace =} \ConstantTok{TRUE}\NormalTok{), }\AttributeTok{sep =} \StringTok{""}\NormalTok{)\}}
\ControlFlowTok{else}\NormalTok{ \{}\FunctionTok{return}\NormalTok{(ans)\}}

\FunctionTok{cat}\NormalTok{(}\StringTok{"}\SpecialCharTok{\textbackslash{}n}\StringTok{"}\NormalTok{)}
\NormalTok{\}}

\CommentTok{\#same thing as before modified for protein}
\FunctionTok{generate\_protein}\NormalTok{(}\DecValTok{66}\NormalTok{,T)}
\end{Highlighting}
\end{Shaded}

\begin{verbatim}
>len66
DFQASEYIPQCKDDQDKPSHMYLQRCHANIYDGRSYDIHMSRWMKKTEDLWGVMKNYPGRLYFENA
\end{verbatim}

Q.4

Generate random protein sequences of length 6 to 13

\begin{Shaded}
\begin{Highlighting}[]
\ControlFlowTok{for}\NormalTok{ (len }\ControlFlowTok{in} \DecValTok{6}\SpecialCharTok{:}\DecValTok{13}\NormalTok{) \{}
 \FunctionTok{cat}\NormalTok{( }\FunctionTok{generate\_protein}\NormalTok{(len, T))}
  \FunctionTok{cat}\NormalTok{(}\StringTok{"}\SpecialCharTok{\textbackslash{}n}\StringTok{"}\NormalTok{)}
  
\NormalTok{\}}
\end{Highlighting}
\end{Shaded}

\begin{verbatim}
>len6
TYWLRR

>len7
MDYWCKL

>len8
AHYYDLMV

>len9
HNGTHHAQQ

>len10
HMTWKCMHFL

>len11
MVYWMSKWDHV

>len12
CIQQWAEPERPY

>len13
MNEKKSRFNTYKF
\end{verbatim}

\begin{Shaded}
\begin{Highlighting}[]
\CommentTok{\# for function allows us to do the same function multiple times without repeating it 8 times}
\end{Highlighting}
\end{Shaded}

Q.5

\subsection{Are our sequences found in
nature?}\label{are-our-sequences-found-in-nature}

\begin{longtable}[]{@{}
  >{\raggedright\arraybackslash}p{(\linewidth - 6\tabcolsep) * \real{0.1857}}
  >{\raggedright\arraybackslash}p{(\linewidth - 6\tabcolsep) * \real{0.3000}}
  >{\raggedright\arraybackslash}p{(\linewidth - 6\tabcolsep) * \real{0.3000}}
  >{\raggedright\arraybackslash}p{(\linewidth - 6\tabcolsep) * \real{0.2143}}@{}}
\toprule\noalign{}
\begin{minipage}[b]{\linewidth}\raggedright
Length (aa)
\end{minipage} & \begin{minipage}[b]{\linewidth}\raggedright
Best hit \% identity
\end{minipage} & \begin{minipage}[b]{\linewidth}\raggedright
Best hit \% coverage
\end{minipage} & \begin{minipage}[b]{\linewidth}\raggedright
Unique? (Y/N)
\end{minipage} \\
\midrule\noalign{}
\endhead
\bottomrule\noalign{}
\endlastfoot
6 & 100 & 100 & N \\
7 & 100 & 100 & N \\
8 & 100 & 88 & Y \\
9 & 88.89 & 100 & Y \\
10 & 90.00 & 100 & Y \\
11 & 88.89 & 82 & Y \\
12 & 88.89 & 75 & Y \\
13 & 88.89 & 69 & Y \\
\end{longtable}

Q.4A

After length 7, the sequences became unique in nature as either identity
or coverage began to drop from 100\%

Q.4B

The greater you increase the number of amino acids in the chain, the
less likely will they be able to line up perfectly with a known
sequence. This comes from the fact that there are 20 different
possibilities of amino acid, thus adding length to the chain now
severely decreases the probability that the entire sequence will line up
with a known sequence, this can be modeled by 20\^{}(length of peptide).
Additionally, the known protein universe itself is limited to what
humans have been able to sequence, although it grows everyday, there are
way more proteins that haven't been sequenced

Q.6

As we can see throughout the class, most people got unique sequences
starting around 7-9 amino acids. Thus it would be best to put the limit
there, at 8. Any shorter, the immune system might accidentally match
sequences that are from the self and thus cause self damage




\end{document}
